\section{维数: 定义和初步性质}\label{sec:dimension-of-affine-varieties}

有了函数域和超越次数的概念, 现在可以给出维数的正式定义.

\begin{Definition}{仿射簇的维数}\leavevmode\label{def:dimension-of-affine-variety}
不可约仿射簇 $X\subseteq\mathbb{A}^{n}$ 的\emph{维数}是它的函数域在 $K$ 上的超越次数:
\[
  \dim(X)=\operatorname{trdeg}_{K}(K(X)).
\]
非空可约仿射簇的\emph{维数}是其各不可约分支维数的最大值.
\end{Definition}

关于维数有许多常用术语. 例如, 一维簇和二维簇通常分别称为\emph{曲线}和\emph{曲面}, $n$ 维簇常称为 \emph{$n$-fold}. 仿射簇 $X\subseteq\mathbb{A}^{n}$ 的\emph{余维数}定义为
\[
  \operatorname{codim}(X)=n-\dim(X),
\]
余维数为一的簇通常称为\emph{超曲面}. 更一般地, 若 $X\subseteq Y\subseteq\mathbb{A}^{n}$ 为仿射簇, 则 $X$ 在 $Y$ 中的\emph{余维数}定义为
\[
  \operatorname{codim}_{Y}(X)=\dim(Y)-\dim(X).
\]
下面在几个具体例子中计算维数.

\begin{Example}{仿射空间 $\mathbb{A}^{n}$ 的维数为 $n$}\leavevmode\label{ex:dimension-of-affine-space}
如 Example~\ref{ex:function-field-of-affine-space} 所见, $\mathbb{A}^{n}$ 的函数域是有理函数域:
\[
  K(\mathbb{A}^{n})=K(x_{1},\ldots,x_{n}).
\]
由于 $x_{1},\ldots,x_{n}\in K(\mathbb{A}^{n})$ 代数独立且生成 $K(\mathbb{A}^{n})$, 它们组成一个超越基. 因此,
\[
  \dim(\mathbb{A}^{n})=\operatorname{trdeg}_{K}(K(\mathbb{A}^{n}))=n.
\]
\end{Example}

\begin{Example}{单位球面的维数为 $2$}\leavevmode\label{ex:dimension-of-unit-sphere}
令 $K=\mathbb{C}$, 并令 $X=\mathcal{V}(x^{2}+y^{2}+z^{2}-1)\subseteq\mathbb{A}^{3}$. 如 Example~\ref{ex:transcendence-basis-of-unit-sphere} 所见, 两个元素
\[
  [x],[y]\in\mathbb{C}(X)
\]
组成一个超越基. 因此,
\[
  \dim(X)=\operatorname{trdeg}_{\mathbb{C}}(\mathbb{C}(X))=2.
\]
更一般地, 可证明仿射簇 $X\subseteq\mathbb{A}^{n}$ 能对某个非零多项式 $f\in K[x_{1},\ldots,x_{n}]$ 写成 $X=\mathcal{V}(f)$, 当且仅当 $X$ 的每个不可约分支都有维数 $n-1$ (Exercise~\ref{ex:characterization-of-hypersurfaces}).
\end{Example}

我们需要借助函数域确保超越次数是良定义的. 不过, 现在基础已经建立, 就能完全用坐标环来讨论维数. 下一个结论说明如何利用 (可能可约的) 簇的坐标环确定其维数的界.

\begin{Proposition}{多项式函数与维数}\leavevmode\label{prop:dimension-bounds-from-coordinate-ring}
设 $X\subseteq\mathbb{A}^{n}$ 为非空仿射簇, 且 $S\subseteq K[X]$.
\begin{enumerate}
  \item\label{part:coordinate-ring-dimension-lower-bound} 若 $S$ 在 $K$ 上代数独立, 则 $|S|\leq\dim(X)$.
  \item\label{part:coordinate-ring-dimension-upper-bound} 若 $K[X]$ 在 $K[S]$ 上整, 则 $\dim(X)\leq|S|$.
\end{enumerate}
\end{Proposition}
\begin{Proof}
证明第一个陈述, 第二个留给 Exercise~\ref{ex:prove-dimension-upper-bound-by-integrality}. 设 $S\subseteq K[X]$ 在 $K$ 上代数独立. 考虑不可约分解 $X=X_{1}\cup\cdots\cup X_{m}$. 将任意函数 $F\in K[X]$ 的定义域限制到分支 $X_{j}$, 得到函数 $F|_{X_{j}}\in K[X_{j}]$. 给定互不相同的函数 $F_{1},\ldots,F_{d}\in S$, 我们断言对至少一个 $j$, 限制
\[
  F_{1}|_{X_{j}},\ldots,F_{d}|_{X_{j}}\in K[X_{j}]
\]
在 $K$ 上代数独立. 为证明这一点, 反设不成立. 那么对每个 $j=1,\ldots,m$, 存在非零多项式 $g_{j}\in K[x_{1},\ldots,x_{d}]$ 使得
\[
  g_{j}(F_{1},\ldots,F_{d})|_{X_{j}}
    =g_{j}(F_{1}|_{X_{j}},\ldots,F_{d}|_{X_{j}})=0.
\]
但 $X$ 是 $X_{1},\ldots,X_{m}$ 的并, 因此乘积 $g_{1}\cdots g_{m}\in K[x_{1},\ldots,x_{d}]$ 给出 $X$ 上一个非平凡的多项式关系:
\[
  g_{1}(F_{1},\ldots,F_{d})\cdots g_{m}(F_{1},\ldots,F_{d})=0\in K[X],
\]
这与 $S$ 在 $K$ 上代数独立的假设矛盾.

因此, 给定互不相同的元素 $F_{1},\ldots,F_{d}\in S$, 至少对一个 $j$, 函数 $F_{1}|_{X_{j}},\ldots,F_{d}|_{X_{j}}$ 在 $K$ 上代数独立. 对这样的 $j$, 得
\[
  \dim(X)\geq\dim(X_{j})=\operatorname{trdeg}_{K}(K(X_{j}))\geq d,
\]
其中第一个不等号和等号由维数的定义得到, 第二个不等号由 Corollary~\ref{cor:bounds-on-transcendence-degree} 第~\ref{part:transcendence-degree-lower-bound} 项得到. 因为这对 $S$ 中任意 $d$ 个不同元素成立, 可知 $|S|\leq\dim(X)$, 如所断言.
\end{Proof}

由 Proposition~\ref{prop:dimension-bounds-from-coordinate-ring} 得到下面的结论: 一个坐标环的每个 Noether 基的大小相同, 且等于该仿射簇的维数.

\begin{Corollary}{Noether 基与维数}\leavevmode\label{cor:noether-basis-counts-dimension}
若 $X\subseteq\mathbb{A}^{n}$ 为非空仿射簇, 且 $\{F_{1},\ldots,F_{d}\}\subseteq K[X]$ 为大小等于 $d$ 的 Noether 基, 则 $d=\dim(X)$.
\end{Corollary}
\begin{Proof}
由 Noether 基的定义, $F_{1},\ldots,F_{d}$ 在 $K$ 上代数独立, 且 $K[X]$ 在 $K[F_{1},\ldots,F_{d}]$ 上整. 因此 Proposition~\ref{prop:dimension-bounds-from-coordinate-ring} 的两部分分别蕴含 $d\leq\dim(X)$ 和 $\dim(X)\leq d$, 所以 $d=\dim(X)$.
\end{Proof}

还可以把维数的正式讨论与 Key Idea~\ref{idea:dimension-as-algebraic-independence} 给出的维数刻画联系起来. 那正是下一结果的第一部分.

\begin{Corollary}{维数的刻画}\leavevmode\label{cor:dimension-and-algebraic-independence}
设 $X\subseteq\mathbb{A}^{n}$ 为非空仿射簇.
\begin{enumerate}
  \item $X$ 的维数是 $K[X]$ 中代数独立函数的最大个数.
  \item\label{part:dimension-as-minimum-integral-generators} $X$ 的维数是满足下列性质的 $F_{1},\ldots,F_{d}\in K[X]$ 的最小个数: $K[X]$ 在 $K[F_{1},\ldots,F_{d}]$ 上整.
\end{enumerate}
\end{Corollary}
\begin{Proof}
证明第一部分, 第二部分留给 Exercise~\ref{ex:prove-dimension-minimum-integral-generators}.

若 $S\subseteq K[X]$ 代数独立, 则 Proposition~\ref{prop:dimension-bounds-from-coordinate-ring} 的第一部分给出 $|S|\leq\dim(X)$. 因此 $K[X]$ 中代数独立函数的最大个数以 $\dim(X)$ 为上界. 另一方面, Noether 正规化告诉我们, $K[X]$ 含有一个 Noether 基, 而由 Corollary~\ref{cor:noether-basis-counts-dimension}, 这个基的大小为 $\dim(X)$. 因而 $K[X]$ 中至少有一个大小为 $\dim(X)$ 的代数独立集. 合起来就得到, $K[X]$ 中代数独立元素的最大个数等于 $\dim(X)$.
\end{Proof}

现在已有几种理解仿射簇维数的方式, 可以借此证明维数满足某些自然的性质. 第一条性质涉及仿射簇之间的包含关系. 由于维数衡量仿射簇的大小, 自然预期 $X\subseteq Y$ 蕴含 $\dim(X)\leq\dim(Y)$. 不那么显然的是, 在不可约仿射簇的情形, 真包含会给出严格不等式. 换句话说, 若 $Y\subseteq\mathbb{A}^{n}$ 是不可约仿射簇, 则真包含于 $Y$ 的仿射簇的维数都严格小于 $Y$ 的维数. 这就是下一个结论.

\begin{Proposition}{维数与包含关系}\leavevmode\label{prop:dimension-monotonicity}
若 $X\subseteq Y$ 为 $\mathbb{A}^{n}$ 中的非空仿射簇, 则 $\dim(X)\leq\dim(Y)$. 若进一步有 $Y$ 不可约且 $X\subsetneq Y$, 则 $\dim(X)<\dim(Y)$.
\end{Proposition}
\begin{Proof}
设 $X,Y\subseteq\mathbb{A}^{n}$ 为非空仿射簇, 且 $X\subseteq Y$. 定义 $d=\dim(X)$. 由 Corollary~\ref{cor:dimension-and-algebraic-independence}, $K[X]$ 中有一个大小为 $d$ 的代数独立子集. 取多项式 $f_{1},\ldots,f_{d}\in K[x_{1},\ldots,x_{n}]$, 使得 $\{f_{1}|_{X},\ldots,f_{d}|_{X}\}$ 为 $K[X]$ 的代数独立子集. 由于 $f_{1},\ldots,f_{d}$ 限制到 $X$ 时不满足任何多项式关系, 限制到较大的集合 $Y$ 时当然也不满足任何多项式关系 (建议读者停下来确认这一点). 因此 $\{f_{1}|_{Y},\ldots,f_{d}|_{Y}\}$ 是 $K[Y]$ 中的代数独立集, 由 Corollary~\ref{cor:dimension-and-algebraic-independence} 得到 $\dim(Y)\geq d=\dim(X)$.

为证明第二个陈述, 进一步设 $Y$ 不可约且 $X\subsetneq Y$. 我们要证明 $\dim(X)<\dim(Y)$. 反设 $\dim(Y)=\dim(X)=d$. 由 $X\subsetneq Y$, 存在 $f\in K[x_{1},\ldots,x_{n}]$ 使得 $f\in\mathcal{I}(X)\setminus\mathcal{I}(Y)$, 换句话说, $f|_{Y}\neq0\in K[Y]$, 但 $f|_{X}=0\in K[X]$. 由于 $K[Y]$ 中代数独立元素的最大个数为 $\dim(Y)=d$, 集合 $\{f_{1}|_{Y},\ldots,f_{d}|_{Y},f|_{Y}\}$ 不可能代数独立 (这里的 $f_{1},\ldots,f_{d}$ 是上一段选取的多项式). 因而可取非零多项式
\[
  g(z)\in K[f_{1}|_{Y},\ldots,f_{d}|_{Y}][z]
\]
使得 $g(f|_{Y})=0$. 不妨设 $g$ 是具有此性质且次数最小的多项式. 逐项写出这一关系:
\[
  \begin{aligned}
    0=g(f|_{Y})={}&g_{k}(f_{1}|_{Y},\ldots,f_{d}|_{Y})(f|_{Y})^{k}+\cdots\\
                &{}+g_{1}(f_{1}|_{Y},\ldots,f_{d}|_{Y})f|_{Y}
                  +g_{0}(f_{1}|_{Y},\ldots,f_{d}|_{Y}),
  \end{aligned}
\]
其中 $g_{0},\ldots,g_{k}\in K[z_{1},\ldots,z_{d}]$. 注意 $g_{0}$ 不可能是零多项式. 否则, 利用 $Y$ 不可约从而 $K[Y]$ 是整环这一假设, 可以约去一个因子 $f|_{Y}$, 得到一个次数更小且在 $f|_{Y}$ 处取值为零的多项式, 与选取 $g$ 时的极小性矛盾. 限制到 $X$ 后, 因为 $f|_{X}$ 是零函数, 所以 $g_{0}$ 给出 $f_{1}|_{X},\ldots,f_{d}|_{X}$ 之间的多项式关系, 与它们是 $K[X]$ 中代数独立元素的假设矛盾. 因此 $\dim(X)<\dim(Y)$, 如所断言.
\end{Proof}

注意, 若不假设 $Y$ 不可约, Proposition~\ref{prop:dimension-monotonicity} 的第二个陈述就不成立. 例如 $x$ 轴 $\mathcal{V}(y)\subseteq\mathbb{A}^{2}$ 真包含于两个坐标轴的并 $\mathcal{V}(xy)\subseteq\mathbb{A}^{2}$, 但这两个簇的维数都是一.

现在有了维数的严格定义, 不难证明它满足 Definition~\ref{def:axioms-for-dimension} 的前两条公理 (见 Exercise~\ref{ex:isomorphic-varieties-have-equal-dimension}, \ref{ex:zero-dimensional-variety-finite} 和~\ref{ex:dimension-of-finite-union}). 所以只剩证明公理~\ref{ax:dimension-of-a-hypersurface-section}, 这就是本章最后一节的内容.

\BookExerciseHeading
\begin{BookExercises}
  \item\label{ex:isomorphic-varieties-have-equal-dimension} 利用本节的结论, 证明两个同构的 (可能可约的) 仿射簇有相同的维数.
  \item\label{ex:zero-dimensional-variety-finite} 设 $X$ 为仿射簇. 证明 $\dim(X)=0$ 当且仅当 $X$ 有限.
  \item\label{ex:dimension-of-finite-union} 设 $X_{1},\ldots,X_{m}\subseteq\mathbb{A}^{n}$ 为 (可能可约的) 仿射簇. 证明
  \[
    \dim(X_{1}\cup\cdots\cup X_{m})=\max\{\dim(X_{1}),\ldots,\dim(X_{m})\}.
  \]
  \item\label{ex:characterization-of-hypersurfaces} 设 $X\subseteq\mathbb{A}^{n}$ 为仿射簇. 证明可对某个非零多项式 $f\in K[x_{1},\ldots,x_{n}]$ 写成 $X=\mathcal{V}(f)$, 当且仅当 $X$ 的每个不可约分支都有维数 $n-1$.
  \item\label{ex:prove-dimension-upper-bound-by-integrality} 证明 Proposition~\ref{prop:dimension-bounds-from-coordinate-ring} 第~\ref{part:coordinate-ring-dimension-upper-bound} 项.
  \item\label{ex:prove-dimension-minimum-integral-generators} 证明 Corollary~\ref{cor:dimension-and-algebraic-independence} 第~\ref{part:dimension-as-minimum-integral-generators} 项.
\end{BookExercises}
